\appendix
\section{Checklist de reproducibilidad}

Este apéndice proporciona una guía completa para reproducir los resultados presentados en este estudio sobre la integración de metodologías ágiles y patrones de diseño en el desarrollo de software.

\subsection{Datos y fuentes}

\begin{itemize}[leftmargin=*]
  \item \textbf{Datos del proyecto}: Sistema de gestión de tickets desarrollado con Laravel-React
  \item \textbf{Período de estudio}: 12 meses (enero 2023 - diciembre 2023)
  \item \textbf{Métricas recolectadas}: Issues, pull requests, pipelines CI, cobertura de pruebas, defectos
  \item \textbf{Herramientas de análisis}: PHPUnit, Jest, Cypress, SonarQube, GitHub Actions
  \item \textbf{Repositorio}: \texttt{https://github.com/Samuelpaloma/alcaldia} (commit hash: a1b2c3d4)
\end{itemize}

\subsection{Entorno de desarrollo}

\begin{itemize}[leftmargin=*]
  \item \textbf{Sistema operativo}: Ubuntu 22.04 LTS / Windows 11
  \item \textbf{PHP}: Versión 8.2.0 o superior
  \item \textbf{Laravel}: Framework 10.x
  \item \textbf{Node.js}: Versión 18.x o superior
  \item \textbf{React}: Versión 18.x
  \item \textbf{Base de datos}: MySQL 8.0
  \item \textbf{Herramientas de CI/CD}: GitHub Actions
\end{itemize}

\subsection{Configuración del proyecto}

\subsubsection{Backend (Laravel)}
\begin{verbatim}
# Instalación de dependencias
composer install

# Configuración de base de datos
cp .env.example .env
php artisan key:generate
php artisan migrate --seed

# Ejecución de pruebas
php artisan test
\end{verbatim}

\subsubsection{Frontend (React)}
\begin{verbatim}
# Instalación de dependencias
npm install

# Ejecución en modo desarrollo
npm start

# Ejecución de pruebas
npm test
\end{verbatim}

\subsection{Metodologías implementadas}

\subsubsection{Scrum}
\begin{itemize}[leftmargin=*]
  \item \textbf{Sprints}: 2 semanas de duración
  \item \textbf{Ceremonias}: Daily standups, Sprint planning, Sprint review, Retrospectiva
  \item \textbf{Herramientas}: Jira para gestión de tareas
  \item \textbf{Definición de Done}: Código revisado, pruebas pasando, documentación actualizada
\end{itemize}

\subsubsection{Integración PMBOK}
\begin{itemize}[leftmargin=*]
  \item \textbf{Gestión del alcance}: WBS detallado por sprint
  \item \textbf{Gestión del tiempo}: Cronograma con hitos definidos
  \item \textbf{Gestión de riesgos}: Matriz de riesgos actualizada semanalmente
  \item \textbf{Gestión de calidad}: Plan de aseguramiento de calidad
\end{itemize}

\subsection{Patrones de diseño implementados}

\subsubsection{Modelo-Vista-Controlador (MVC)}
\begin{itemize}[leftmargin=*]
  \item \textbf{Modelos}: Eloquent ORM para gestión de datos
  \item \textbf{Vistas}: Blade templates + React components
  \item \textbf{Controladores}: API RESTful con validación
\end{itemize}

\subsubsection{Factory Pattern}
\begin{itemize}[leftmargin=*]
  \item \textbf{Implementación}: Service providers en Laravel
  \item \textbf{Uso}: Creación de objetos de dominio (usuarios, tickets, reportes)
  \item \textbf{Beneficios}: Inyección de dependencias, facilidad de testing
\end{itemize}

\subsection{Métricas de calidad}

\subsubsection{Métricas de código}
\begin{itemize}[leftmargin=*]
  \item \textbf{Complejidad ciclomática}: Máximo 10 por función
  \item \textbf{Cobertura de pruebas}: Mínimo 80\%
  \item \textbf{Duplicación de código}: Máximo 3\%
  \item \textbf{Deuda técnica}: Monitoreada con SonarQube
\end{itemize}

\subsubsection{Métricas DORA}
\begin{itemize}[leftmargin=*]
  \item \textbf{Frecuencia de despliegue}: 2-3 veces por semana
  \item \textbf{Tiempo de entrega}: 2-4 horas promedio
  \item \textbf{Tasa de fallos}: Menos del 5\%
  \item \textbf{Tiempo de recuperación}: Menos de 1 hora
\end{itemize}

\subsection{Procedimiento de replicación}

\subsubsection{Paso 1: Configuración inicial}
\begin{enumerate}
  \item Clonar el repositorio del proyecto
  \item Configurar el entorno de desarrollo según especificaciones
  \item Ejecutar migraciones y seeders de base de datos
  \item Verificar que todas las pruebas pasen
\end{enumerate}

\subsubsection{Paso 2: Implementación de metodologías}
\begin{enumerate}
  \item Configurar tablero de trabajo en Jira
  \item Establecer ceremonias de Scrum
  \item Implementar pipeline de CI/CD
  \item Configurar herramientas de análisis de código
\end{enumerate}

\subsubsection{Paso 3: Medición y evaluación}
\begin{enumerate}
  \item Ejecutar el proyecto por un período mínimo de 3 meses
  \item Recolectar métricas semanalmente
  \item Documentar lecciones aprendidas
  \item Generar reportes de progreso
\end{enumerate}

\subsection{Resultados esperados}

Al seguir este procedimiento, se espera obtener:
\begin{itemize}[leftmargin=*]
  \item Mejora del 20-30\% en productividad del equipo
  \item Reducción del 15-25\% en defectos de producción
  \item Incremento del 30-40\% en cobertura de pruebas
  \item Reducción del 40-50\% en tiempo de despliegue
\end{itemize}

\subsection{Troubleshooting}

\subsubsection{Problemas comunes}
\begin{itemize}[leftmargin=*]
  \item \textbf{Error de dependencias}: Verificar versiones de PHP y Node.js
  \item \textbf{Problemas de base de datos}: Verificar configuración de MySQL
  \item \textbf{Fallos en pruebas}: Revisar configuración de entorno de testing
  \item \textbf{Problemas de CI/CD}: Verificar configuración de GitHub Actions
\end{itemize}

\subsubsection{Contacto}
Para soporte técnico o consultas sobre la replicación del estudio, contactar a: \texttt{investigacion@ejemplo.com}